% --- section_1.tex ---
\begin{tabularx}{\columnwidth}{XX}
  \multicolumn{2}{l}{\textbf{基本概念}} \\
  瞬时速度 & $v=\lim\limits_{t\to0} \frac{\Delta x}{\Delta t}$ \\
  平均速度 & $v=\frac{\Delta x}{\Delta t}$ \\
  瞬时速率 & $v=$瞬时速度大小 \\
  平均速率 & $v=\frac{\text{路程}}{\Delta t}$ \\
  平均加速度 & $a=\frac{\Delta v}{\Delta t}$ \\
  \multicolumn{2}{l}{\textbf{匀变速直线运动}} \\
  位移 & $x=v_0t+\frac{1}{2}at^2$ \\
   & $x=\frac{v_0+v}{2}t$ \\
  速度 & $v=v_0+at=v_tt-at$ \\
  速度平方 & $v^2=v_0^2+2ax$ \\
  中间位移速度 & $v=\sqrt{\frac{x_1^2+x_2^2}2}$ \\
  中间时刻速度 & $v=\frac{x_1+x_2}{2}$ \\
  相邻相等时间间隔位移差 & $\Delta x=aT^2$ \\
  第M段与第N段Ts时间内位移差 & $(M-N)aT^2$ \\
  \multicolumn{2}{l}{\textbf{从静止开始加速}} \\
  相等时间每段位移比 & $1:3:5:\cdots:2n-1$ \\
  相等时间总计位移比 & $1:4:9:\cdots:n^2$ \\
  相等位移间隔每段用时比 & $1:\sqrt2-1:\cdots:\sqrt n-\sqrt{n-1}$ \\
  相等位移间隔总计用时比 & $1:\sqrt2:\cdots:\sqrt n$ \\
\end{tabularx}
\medskip

% --- section_2.tex ---
\begin{tabularx}{\columnwidth}{XX}
  水平位移 & $x=v_0 t$ \\
  竖直位移 & $y=\frac12 a t^2$ \\
  竖直速度 & $v_y=a t=\sqrt{\frac{2y}{a}}=v_0\tan\theta$ \\
  合速度 & $v=\sqrt{v_y^2+v_0^2}=\frac{v_0}{\cos\theta}$ \\
  轨迹抛物线方程 & $y=\frac{ax^2}{2v_0^2}$ \\
  \multicolumn{2}{l}{\textbf{平抛运动推论}} \\
  速度偏角 & $\tan\theta=\frac{v_y}{v_0}=\frac{gt}{v_0}=\frac{2y}{x}$ \\
  斜面抛回斜面 & $\tan\theta=2\tan\text{斜面倾角}$ \\
  垂直打到斜面上 & $\tan\theta \tan\alpha = 1 $ \\
\end{tabularx}
\medskip

% --- section_3.tex ---
\begin{tabularx}{\columnwidth}{XX}
  向心力 & $F_{\text{向}}=m\frac{v^2}{r}$ \\
  向心力 & $F_{\text{向}}=ma_{\text{向}}$ \\
  线速度 & $v=\omega r$ \\
  角速度 & $\omega=\frac{2\pi}{T}$ \\
\end{tabularx}
\medskip

% --- section_4.tex ---
\begin{tabularx}{\columnwidth}{XX}
  球体万有引力 & $ F_{\text{万}}=\left\{
    \begin{cases}
      \frac{GMm}{r^2}  & r\ge R \\
      \frac{GMmr}{R^3} & r<R    \\
    \end{cases}
  \right.$ \\
  球体万有引力势能 & $E_p=-\frac{GMm}{r},r\ge R$ \\
  向心力简化式 & $GM=v^2r$ \\
   & $GM=\omega^2r^3$ \\
   & $GM=ar^2$ \\
   & $GM=4\pi^2\frac{r^3}{T^2}$ \\
  密度公式 & $\rho=\frac{3\pi}{GT^2}$ \\
  \multicolumn{2}{l}{\textbf{开普勒三大定律}} \\
  周期定律 & $\frac{a^3}{T^2}=k=\frac{GM}{4\pi^2}$ \\
  \multicolumn{2}{l}{\textbf{黄金代换}} \\
  黄金代换 & $GM=gR^2$ \\
  \multicolumn{2}{l}{\textbf{宇宙速度}} \\
  第一宇宙速度 & $v=\sqrt{\frac{GM}{R}}=\sqrt{Rg}$ \\
  \multicolumn{2}{l}{\textbf{追及问题}} \\
  同向（最近到最远） & $t=\frac12\frac{T_{\text{大}}T_{\text{小}}}{T_{\text{大}}-T_{\text{小}}}$ \\
  同向（最近到最近） & $t=\frac{T_{\text{大}}T_{\text{小}}}{T_{\text{大}}-T_{\text{小}}}$ \\
  反向（最近到最远） & $t=\frac12\frac{T_{\text{大}}T_{\text{小}}}{T_{\text{大}}+T_{\text{小}}}$ \\
  反向（最近到最近） & $t=\frac{T_{\text{大}}T_{\text{小}}}{T_{\text{大}}+T_{\text{小}}}$ \\
  \multicolumn{2}{l}{\textbf{双星问题}} \\
  质量半径关系 & $m_1r_1=m_2r_2$ \\
  质量角速度关系 & $G(m_1+m_2)=\omega^2(r_1+r_2)^3$ \\
\end{tabularx}
\medskip

% --- section_5.tex ---
\begin{tabularx}{\columnwidth}{XX}
  重力 & $G=mg$ \\
  两极重力加速度 & $g_{\text{两极}}=\frac{GM}{R^2}$ \\
  赤道重力加速度 & $g_{\text{赤道}}=\frac{GM}{R^2}-\omega^2R$ \\
  滑动摩擦力 & $f=\mu N$ \\
  最大静摩擦力 & $f_{\max}=\mu N$ \\
  弹簧弹力 & $F=kx$ \\
  共点力平衡 & $\sum F_x=0,\quad \sum F_y=0$ \\
\end{tabularx}
\medskip

% --- section_6.tex ---
\begin{tabularx}{\columnwidth}{XX}
  \multicolumn{2}{l}{\textbf{牛顿运动定律}} \\
  牛顿第二定律 & $\sum F=ma$ \\
  牛顿第三定律 & $F_{AB}=-F_{BA}$ \\
  重力 & $G=mg$ \\
  滑动摩擦力 & $f=\mu N$ \\
  最大静摩擦力 & $f_{\max}\approx \mu_s N$ \\
  胡克定律 & $F=kx$ \\
  \multicolumn{2}{l}{\textbf{空气阻力}} \\
  空气阻力 & $f=kv$ \\
\end{tabularx}
\medskip

% --- section_7.tex ---
\begin{tabularx}{\columnwidth}{XX}
  动能定理 & $W_{\text{合力}}=\Delta E_k$ \\
  动能 & $E_k=\frac12mv^2$ \\
  重力做功 & $W_G=mg(h_1-h_2)$ \\
  弹力做功 & $W_N=0$（弹力始终垂直速度时） \\
\end{tabularx}
\medskip

% --- section_8.tex ---
\begin{tabularx}{\columnwidth}{XX}
  功 & $W=Fs\cos\theta$ \\
  合力做功 & $W_{\text{合}}=\Delta E_k$ \\
  动能 & $E_k=\frac12mv^2$ \\
  重力势能 & $E_p=mgh$ \\
  弹性势能 & $E_p=\frac12kx^2$ \\
  重力做功与重力势能 & $W_G=-\Delta E_{pG}=mg(h_1-h_2)$ \\
  弹力做功与弹性势能 & $W_{\text{弹}}=-\Delta E_{p\text{弹}}$ \\
  机械能 & $E=E_k+E_p$ \\
  机械能守恒 & $E_{k1}+E_{p1}=E_{k2}+E_{p2}$ \\
  非保守力做功 & $W_{\text{非保}}=\Delta E_{\text{机}}$ \\
\end{tabularx}
\medskip

% --- section_9.tex ---
\begin{tabularx}{\columnwidth}{XX}
  动量定理 & $\vec I=\Delta\vec p$ \\
  恒力冲量 & $\vec I=\vec F t$ \\
  平均力冲量 & $\vec I=\vec F_{\text{平均}}\Delta t$ \\
  系统动量定理 & $\vec I_{\text{外}}=\Delta\vec P_{\text{系统}}$ \\
\end{tabularx}
\medskip

% --- section_30.tex ---
\begin{tabularx}{\columnwidth}{XX}
  动量守恒 & $m_1v_1+m_2v_2=m_1v_1'+m_2v_2'$ \\
  恢复系数 & $e=\frac{\text{分离相对速度}}{\text{接近相对速度}}$ \\
  一维恢复系数 & $v_2'-v_1'=e(v_1-v_2)$ \\
  约化质量 & $\mu=\frac{m_1m_2}{m_1+m_2}$ \\
  碰撞损失机械能 & $\Delta E=\frac12\mu(1-e^2)(v_1-v_2)^2$ \\
\end{tabularx}
\medskip

% --- section_24.tex ---
\begin{tabularx}{\columnwidth}{XX}
  回复力 & $F=-kx$ \\
  简谐振动方程 & $x(t)=A\sin(\omega t+\varphi_0)$ \\
  速度与加速度 & $v=\omega A\cos(\omega t+\varphi_0)$，$a=-\omega^2x$ \\
  周期频率关系 & $T=\frac{2\pi}{\omega}$，$f=\frac{1}{T}$ \\
  弹簧振子周期 & $T=2\pi\sqrt{\frac{m}{k}}$ \\
  单摆回复力 & $F=-\frac{mg}{l}x$ \\
  单摆周期 & $T=2\pi\sqrt{\frac{l}{g}}$ \\
  水平弹簧振子能量 & $E=\frac{1}{2}kx^2+\frac{1}{2}mv^2=\frac12kA^2$ \\
\end{tabularx}
\medskip

% --- section_25.tex ---
\begin{tabularx}{\columnwidth}{XX}
  波动通式 & $y=A\sin(\omega t\mp kx+\varphi_0)$，$k=\frac{2\pi}{\lambda}$ \\
  任意点振动方程 & $y=A\sin(\omega (t-\frac{x}{v})+\varphi_0)$ \\
  波速波长频率关系 & $v=\lambda f=\frac{\lambda}{T}$ \\
  相位差 & $\Delta\varphi=\frac{2\pi}{\lambda}\Delta x=2\pi\frac{\Delta t}{T}$ \\
  \multicolumn{2}{l}{\textbf{两波源在该点起振方向相同}} \\
  干涉加强条件 & $\Delta x=n\lambda$ \\
  干涉减弱条件 & $\Delta x=n\lambda+\frac12\lambda$ \\
  \multicolumn{2}{l}{\textbf{两波源在该点起振方向相反}} \\
  干涉加强条件 & $\Delta x=n\lambda+\frac12\lambda$ \\
  干涉减弱条件 & $\Delta x=n\lambda$ \\
  衍射条件 & 障碍尺寸越小、波长越大，衍射越明显。 \\
  驻波条件 & 两列同频同振幅机械波相向传播 \\
  多普勒效应 & 相互靠近频率变高；相互远离频率变低。 \\
\end{tabularx}
\medskip

% --- section_10.tex ---
\begin{tabularx}{\columnwidth}{XX}
  库仑定律 & $F=k\dfrac{Qq}{r^2}$ \\
  电场强度定义 & $E=\dfrac{F}{q}$ \\
  点电荷场强 & $E=k\dfrac{Q}{r^2}$ \\
  电场叠加 & $\vec E=\vec E_1+\vec E_2+\cdots$ \\
  匀强电场 & $E=\text{恒量}$ \\
\end{tabularx}
\medskip

% --- section_31.tex ---
\begin{tabularx}{\columnwidth}{XX}
  电势定义 & $\varphi=\dfrac{E_p}{q}$ \\
  电势差定义 & $U_{AB}=\varphi_A-\varphi_B$ \\
  电场力做功 & $W_{AB}=qU_{AB}$ \\
  电势能变化 & $\Delta E_p=-W_{\text{电}}$ \\
  匀强电场电势差 & $U=Ed$ \\
  电场力做功 & $W=qEd\cos\theta$ \\
\end{tabularx}
\medskip

% --- section_32.tex ---
\begin{tabularx}{\columnwidth}{XX}
  电容定义 & $C=\dfrac{Q}{U}$ \\
  平行板电容器 & $C=\varepsilon_0\varepsilon_r\dfrac{S}{d}$ \\
  电容器能量 & $E_C=\dfrac12CU^2=\dfrac{Q^2}{2C}=\dfrac12QU$ \\
  加速电场 & $qU=\dfrac12mv^2$ \\
  偏转电场加速度 & $a=\dfrac{qE}{m}=\dfrac{qU}{md}$ \\
  偏转位移 & $y=\dfrac12at^2$ \\
\end{tabularx}
\medskip

% --- section_11.tex ---
\begin{tabularx}{\columnwidth}{XX}
  电流定义式 & $I=\dfrac{q}{t}$ \\
  电流微观表达式 & $I=nqSv$ \\
  电阻定义式 & $R=\dfrac{U}{I}$ \\
  电阻决定式 & $R=\rho\dfrac{l}{S}$ \\
  串联电阻 & $R=R_1+R_2+\cdots$ \\
  并联电阻 & $\dfrac1R=\dfrac1{R_1}+\dfrac1{R_2}+\cdots$ \\
\end{tabularx}
\medskip

% --- section_33.tex ---
\begin{tabularx}{\columnwidth}{XX}
  电动势 & $E=\dfrac{W_{\text{非静电力}}}{q}$ \\
  闭合电路欧姆定律 & $E=I(R+r)$ \\
  路端电压 & $U=E-Ir$ \\
  纯电阻功率 & $P=UI=I^2R=\dfrac{U^2}{R}$ \\
  电源总功率 & $P_{\text{总}}=EI$ \\
  电源输出功率 & $P_{\text{出}}=UI$ \\
  电源效率 & $\eta=\dfrac{P_{\text{出}}}{P_{\text{总}}}=\dfrac{U}{E}$ \\
\end{tabularx}
\medskip

% --- section_34.tex ---
\begin{tabularx}{\columnwidth}{XX}
  电流表改装分流电阻 & $R=\dfrac{I_gR_g}{I-I_g}$ \\
  电压表改装串联电阻 & $R=\dfrac{U}{I_g}-R_g$ \\
  欧姆表回路 & $E=i(R_{\text{内}}+R_x)$ \\
  欧姆表中值电阻 & $R_{\text{中}}=R_{\text{内}}$ \\
  电桥平衡条件 & $\dfrac{R_1}{R_2}=\dfrac{R_3}{R_x}$ \\
\end{tabularx}
\medskip

% --- section_12.tex ---
\begin{tabularx}{\columnwidth}{XX}
  磁感应强度 & $B=\dfrac{F}{Il}$ \\
  安培力 & $F=BIl\sin\theta$ \\
  通电导线受力方向 & 左手定则 \\
  直导线磁场方向 & 安培定则 \\
  螺线管磁场方向 & 右手螺旋定则 \\
\end{tabularx}
\medskip

% --- section_35.tex ---
\begin{tabularx}{\columnwidth}{XX}
  洛伦兹力 & $F=qvB\sin\theta$ \\
  向心力关系 & $qvB=m\dfrac{v^2}{R}$ \\
  半径 & $R=\dfrac{mv}{qB}$ \\
  周期 & $T=\dfrac{2\pi m}{qB}$ \\
  角速度 & $\omega=\dfrac{qB}{m}$ \\
  运动时间 & $t=\dfrac{\theta m}{qB}=\dfrac{\theta R}{v}$ \\
\end{tabularx}
\medskip

% --- section_36.tex ---
\begin{tabularx}{\columnwidth}{XX}
  速度选择器 & $qE=qvB$ \\
  选择速度 & $v=\dfrac{E}{B}$ \\
  加速电场 & $qU=\dfrac12mv^2$ \\
  质谱仪半径 & $R=\dfrac1B\sqrt{\dfrac{2mU}{q}}$ \\
  霍尔电压 & $U_H=\dfrac1{nq}\dfrac{IB}{d}$ \\
  回旋加速器最大速度 & $v_{\max}=\dfrac{qBR}{m}$ \\
  回旋加速器最大动能 & $E_{k\max}=\dfrac{q^2B^2R^2}{2m}$ \\
\end{tabularx}
\medskip

% --- section_13.tex ---
\begin{tabularx}{\columnwidth}{XX}
  磁通量 & $\Phi=BS\cos\theta$ \\
  法拉第电磁感应定律 & $E=n\left|\dfrac{\Delta\Phi}{\Delta t}\right|$ \\
  感应电流 & $I=\dfrac{E}{R+r}$ \\
  回路电荷量 & $q=\dfrac{n\Delta\Phi}{R+r}$ \\
  导体棒切割 & $E=Blv$ \\
\end{tabularx}
\medskip

% --- section_37.tex ---
\begin{tabularx}{\columnwidth}{XX}
  动生电动势 & $E=BLv$ \\
  回路电流 & $I=\dfrac{E}{R}$ \\
  安培力 & $F_A=BIL=\dfrac{B^2L^2v}{R}$ \\
  安培力冲量 & $I_A=BLq$ \\
  感应电流功率 & $P=EI=I^2R$ \\
  恒力最终速度 & $v=\dfrac{FR}{B^2L^2}$ \\
\end{tabularx}
\medskip

% --- section_38.tex ---
\begin{tabularx}{\columnwidth}{XX}
  磁通量 & $\Phi=BS\cos\theta$ \\
  线圈转动电动势 & $e=nBS\omega\sin\theta$ \\
  最大电动势 & $E_m=nBS\omega$ \\
  从中性面计时 & $e=E_m\sin\omega t$ \\
  周期与角速度 & $T=\dfrac{2\pi}{\omega}$ \\
  双棒动量关系 & $m_1v_1+m_2v_2=\text{常量}$（无外力且安培力等大反向） \\
\end{tabularx}
\medskip

% --- section_14.tex ---
\begin{tabularx}{\columnwidth}{XX}
  双缝干涉条纹间距 & $\Delta x=\frac{L}{d}\lambda$ \\
  相干加强条件 & $\Delta r=k\lambda\quad(k=0,1,2,\cdots)$ \\
  相干减弱条件 & $\Delta r=(k+\frac12)\lambda$ \\
  光程 & $s=nr$ \\
  洛埃德镜条纹间距 & $\Delta x=\frac{L}{d}\lambda$ \\
  薄膜干涉光程差 & $\Delta=2nd+\delta$（近似垂直入射） \\
  半波损失 & 从光疏到光密反射时 $\delta=\frac{\lambda}{2}$ \\
  劈尖厚度 & $d=x\theta$（小角度，$\theta$ 用弧度） \\
  劈尖条纹间距 & $\Delta x=\frac{\lambda}{2n\theta}$ \\
\end{tabularx}
\medskip

% --- section_15.tex ---
\begin{tabularx}{\columnwidth}{XX}
  反射定律 & $i=r$ \\
  折射定律 & $n_1\sin i=n_2\sin r$ \\
  折射率定义 & $n=\frac{c}{v}$ \\
  相对折射率 & $n_{21}=\frac{n_2}{n_1}=\frac{\sin i}{\sin r}$ \\
  全反射临界角 & $\sin C=\frac{n_2}{n_1}\quad(n_1>n_2)$ \\
  透镜成像公式 & $\frac{1}{f}=\frac{1}{u}+\frac{1}{v}$ \\
  放大率 & $m=\frac{h'}{h}=\frac{v}{u}$ \\
\end{tabularx}
\medskip

% --- section_16.tex ---
\begin{tabularx}{\columnwidth}{XX}
  热力学温度 & $T=273+t_{\text{摄氏度}}$ \\
  微粒数 & $N=nN_A$ \\
  物质的量 & $n=\dfrac{m}{M}$ \\
  单个分子平均占有体积 & $V_{\text{单个}}=\dfrac{V}{N}$ \\
  固体、液体分子直径估算 & $d\sim \sqrt[3]{\frac{6V_{\text{单个}}}{\pi}}$ \\
  气体分子平均间距估算 & $a\sim \sqrt[3]{V_{\text{单个}}}$ \\
  分子平均动能 & $E=\dfrac{3}{2}kT$ \\
  分子平均速率 & $v=\sqrt{\dfrac{3kT}{m}}$ \\
\end{tabularx}
\medskip

% --- section_17.tex ---
\begin{tabularx}{\columnwidth}{XX}
  等温变化 & $pV=C$ \\
  等容变化 & $p=CT$ \\
  等容变化比值形式 & $\frac{p}{T}=\frac{\Delta p}{\Delta T}$ \\
  等压变化 & $V=CT$ \\
  等压变化比值形式 & $\frac{V}{T}=\frac{\Delta V}{\Delta T}$ \\
  理想气体状态方程 & $pV=nRT$ \\
  液体上下压强 & $p_{\text{下}}-p_{\text{上}}=\rho g(h_{\text{下}}-h_{\text{上}})$ \\
  有质量活塞上下压强 & $p_{\text{下}}-p_{\text{上}}=\frac{Mg}{S}$ \\
  水平液柱与活塞压强 & $p_{\text{左}}=p_{\text{右}}$ \\
  竖直液柱加速度 & $p_{\text{下}}-p_{\text{上}}=\rho g_{\text{等效}}h$ \\
  竖直活塞加速度 & $p_{\text{下}}-p_{\text{上}}=\frac{M}{S}g_{\text{等效}}$ \\
\end{tabularx}
\medskip

% --- section_18.tex ---
\begin{tabularx}{\columnwidth}{XX}
  热力学第一定律 & $\Delta U=Q+W$ \\
  等容过程 & $W=0$ \\
  等压过程气体对外做功 & $W_{\text{气}}=p\Delta V$ \\
  压强线性变化做功 & $W_{\text{气}}=\frac{p_1+p_2}{2}\Delta V$ \\
  等温过程 & $\Delta U=0$ \\
  绝热过程 & $Q=0$ \\
\end{tabularx}
\medskip

% --- section_19.tex ---
\begin{tabularx}{\columnwidth}{XX}
  能量守恒定律 & 能量既不会凭空产生，也不会凭空消失 \\
  第一类永动机 & 违反了能量守恒定律，本质上要求无中生有地产生能量 \\
  第二类永动机 & 违反了热力学第二定律，本质上要求热量100\% 转化为有用功 \\
\end{tabularx}
\medskip

% --- section_26.tex ---
\begin{tabularx}{\columnwidth}{XX}
  能量子关系 & $E=h\nu$ \\
  光子频率与波长 & $\nu=\frac{c}{\lambda}$ \\
  光子能量表达式 & $E=\frac{hc}{\lambda}$ \\
  振子能量量子化 & $E_n=nh\nu\quad(n=0,1,2,\cdots)$ \\
  峰值波长规律 & $\lambda_m T=\text{常量}$ \\
  黑体辐射规律（定性） & 温度越高，总辐射越强，峰值波长越短 \\
\end{tabularx}
\medskip

% --- section_20.tex ---
\begin{tabularx}{\columnwidth}{XX}
  光子能量 & $E=h\nu=\frac{hc}{\lambda}$ \\
  光子动量 & $p=\frac{h}{\lambda}=\frac{E}{c}$ \\
  德布罗意波长 & $\lambda=\frac{h}{p}$ \\
  德布罗意关系（非相对论） & $\lambda=\frac{h}{mv}$ \\
\end{tabularx}
\medskip

% --- section_21.tex ---
\begin{tabularx}{\columnwidth}{XX}
  爱因斯坦光电方程 & $h\nu=W_0+E_{k\max}$ \\
  最大初动能（电压表示） & $E_{k\max}=eU_c$ \\
  截止频率 & \(\nu_0=\frac{W_0}{h}\) \\
  逸出功（波长表示） & \(W_0=\frac{hc}{\lambda_0}\) \\
  最大初动能--频率图像 & \(E_{k\max}=h\nu-W_0=h(\nu-\nu_0)\) \\
  遏止电压--频率图像 & \(U_c=\frac{h}{e}\nu-\frac{W_0}{e}=\frac{h}{e}(\nu-\nu_0)\) \\
\end{tabularx}
\medskip

% --- section_22.tex ---
\begin{tabularx}{\columnwidth}{XX}
  玻尔能级公式（氢原子） & $E_n=-\frac{13.6\ \mathrm{eV}}{n^2}$ \\
  轨道半径（氢原子） & $r_n=\frac{r_0}{n^2}$ \\
  频率条件 & $h\nu=|E_m-E_n|$ \\
  最大谱线数（从第 $n$ 能级向下） & $N=\frac{n(n-1)}{2}$ \\
\end{tabularx}
\medskip

% --- section_23.tex ---
\begin{tabularx}{\columnwidth}{XX}
  衰变规律 & $N=N_0\left(\frac12\right)^{t/T}$ \\
  半衰期定义 & $T$：核数减半所需时间 \\
  质量亏损 & $\Delta m=m_{\text{初}}-m_{\text{末}}$ \\
  质能关系 & $E=\Delta mc^2$ \\
  原子质量单位换算 & $1\ \mathrm{u}c^2\approx 931.5\ \mathrm{MeV}$ \\
  结合能法放能 & $Q=\sum E_{b,\text{生成物}}-\sum E_{b,\text{反应物}}$ \\
  比结合能 & $\bar E_b=\frac{E_b}{A}$ \\
\end{tabularx}
\medskip

% --- section_23.tex ---
\begin{tabularx}{\columnwidth}{|c|X|}
  \hline
  反应类型 & 反应方程式 \\
  \hline
  发现质子 & \ce{^4_2He + ^{14}_7N -> ^1_1H + ^{17}_8O} \\
  \hline
  发现中子 & \ce{^4_2He + ^{9}_4Be -> ^1_0n + ^{12}_6C} \\
  \hline
  发现正电子 & \ce{^4_2He + ^{27}_{13}Al -> ^{30}_{15}P + ^1_0n} \\
  \hline
  $\alpha$ 衰变 & \ce{^{238}_{92}U -> ^{234}_{90}Th + ^4_2He} \\
  \hline
  $\beta^-$ 衰变本质 & \ce{^1_0n -> ^1_1p + ^0_{-1}e} \\
  \hline
  核裂变 & \ce{^{235}_{92}U + ^1_0n -> ^{92}_{36}Kr + ^{141}_{56}Ba + 3^1_0n} \\
  \hline
  核聚变 & \ce{^2_1H + ^3_1H -> ^4_2He + ^1_0n} \\
  \hline
\end{tabularx}
\medskip

